%=====================================================================================================
\section{Introduction}
Counting and tracking people in public spaces is a mature computer-vision task, yet a deployed crowd-analytics
system fails in ways a benchmark score does not show. The detector is adapted to one camera with labels nobody
checked by hand; the camera it is later pointed at looks different; a frame-level confidence interval is reported as
if neighbouring frames were independent; and, increasingly, a vision-language model (VLM) is allowed to answer
questions about the numbers and to change the system's configuration. Each of these is a reliability question, and
each has strong components to build on---real-time
detectors~\cite{redmon2016yolo,wang2024yolov10,jocher2026yolo26,zhao2024rtdetr,robinson2026rfdetr},
tracking-by-detection~\cite{bewley2016sort,zhang2022bytetrack,cao2023ocsort}, open
VLMs~\cite{bai2025qwen3vl}---but they are rarely studied together, under one protocol, on the same footage.

This paper studies the reliability of such a pipeline end to end. A teacher ensemble pseudo-labels an unlabelled
5-minute plaza video (\texttt{aa1}); two compact students (CNN \yolo{} and transformer \rfdetr{}) are fine-tuned on
it; the selected student drives an FFmpeg-based edge runtime with ByteTrack, deterministic analytics and a
Qwen3-VL-2B agent that writes scene reports and reconfigures the analytics through tool calls. We then ask where it
fails: under which labels a model looks better, what happens on cameras it was not adapted to, how much the
statistics change when frames are treated as correlated, and what a small VLM does with numbers, with instructions
it reads in the image, and with requests it should refuse. The system targets an NVIDIA Jetson
Orin~\cite{jetsonorin}, but every measurement here is from a desktop GPU; on-device validation is left to future
work.

Most components are established; the contributions are the measurements and the design decisions they force:
\begin{itemize}
\item \textbf{When pseudo-label scene adaptation hurts, and a cheap fix.} Fine-tuning \yolo{} on one camera raises
in-domain AP but lowers it on every other camera we tested and on human-labelled COCO people, below its own
pretrained weights. A size-blind label threshold that turns small, far people into background explains part of this
and forgetting the rest. Of five mitigations across four scenes, weight-space ensembling of pretrained and fine-tuned
weights works best and costs nothing; chosen with a shift-aware selection that never sees the test camera, it keeps
most of the in-domain gain and the pretrained robustness.
\item \textbf{An evaluation protocol for correlated video} with explicit units, circular-block bootstrap intervals
and corrected bootstrap $p$-values, label-bias and GT-linking checks, and an independent re-implementation of every
metric; a resolution-matched comparison shows that the CNN student's in-domain lead was an input-size effect.
\item \textbf{VLM reliability measured, then designed for.} A 2B VLM copies counts and rubber-stamps fabricated
alerts, so counts and alert decisions stay deterministic; tool routing is separated from visual answering; and a
300-case adversarial suite (ambiguous, conflicting, unauthorised, malformed, alert-disabling, injected in the request,
in the frame and in tool results) compares the model alone, the original guardrails and a trusted policy layer with
role-based access, validation, an intent gate, confirmation, rate limits and a hash-chained audit log.
\item \textbf{A desktop-validated, Jetson-oriented runtime} that reproduces the offline tracking accuracy when
streamed through the deployed decoder, ONNX detector and online tracker.
\item \textbf{A reproducibility package}: pinned environment, checksums, per-table commands, and a blind
human-annotation protocol and tool for an independent ground truth.
\end{itemize}
